% ==============================================================================
% 3.6 PRUEBAS (flujo de trabajo PUDS dentro del Capítulo 3)
% Las pruebas sobre el producto terminado y los mecanismos de prueba se desarrollan en el Capítulo 9.
% [INSTRUCCIÓN]: los apartados 3.6.1 a 3.6.6 siguen las actividades del flujo de trabajo de Pruebas del
% Proceso Unificado; ajustarlos si el PUDS de la cátedra usa otros títulos. Las cifras son las medidas el
% 18-nov-2026 (informe de pruebas del repositorio) y se actualizan en el Capítulo 9.
% ==============================================================================
\subsection{3.6. Pruebas}\label{sec:v2_flujo_pruebas}

Este flujo de trabajo del PUDS verifica que el sistema implementado cumpla los requisitos especificados en la sección 3.2. Las pruebas sobre el producto terminado y los mecanismos de prueba utilizados se desarrollan en el Capítulo~\hyperref[cap:pruebas]{9}; aquí se describe cómo se planificaron, diseñaron, implementaron y evaluaron durante el desarrollo.

\subsubsection{3.6.1. Planificar las pruebas}

El objetivo de las pruebas es comprobar que cada uno de los 26 casos de uso funciona como lo describen sus flujos y que el sistema cumple los requisitos no funcionales de rendimiento (pantallas en 3 s o menos y cálculo del puntaje en 5 s o menos). La planificación define cuatro decisiones:

\begin{itemize}
    \item \textbf{Niveles de prueba:} unidad (cada servicio y componente), integración (el caso de uso completo contra una base de datos real), sistema (la API y el panel web en funcionamiento) y rendimiento (carga y tiempos de respuesta).
    \item \textbf{Alcance por parte del sistema:} el \textit{backend}, el panel web, la aplicación móvil y el módulo de inteligencia artificial tienen pruebas automatizadas propias; la API completa y el panel web se prueban además de punta a punta.
    \item \textbf{Entorno:} PostgreSQL real levantado en un contenedor (Testcontainers) en las pruebas de integración; el \textit{backend} en perfil de desarrollo, con datos de prueba, en las pruebas de sistema.
    \item \textbf{Criterios de aceptación:} todas las pruebas automatizadas pasan en la integración continua; cada caso de uso prueba su camino correcto, sus errores de negocio y el permiso por rol; y se cumplen los umbrales de rendimiento.
\end{itemize}

Como el tiempo del proyecto es corto, se priorizaron los casos de uso de prioridad alta (CU-02, CU-03 y CU-05) para el análisis más detallado de la sección 9.3, con técnicas de caja blanca y de caja negra.

\subsubsection{3.6.2. Diseñar las pruebas}

Los casos de prueba se derivan de los casos de uso, con la misma estructura de la especificación de cada uno:

\begin{itemize}
    \item \textbf{Flujo principal:} un caso de prueba del camino correcto.
    \item \textbf{Flujos alternativos y excepciones:} un caso por cada error de negocio, con el código de error esperado (por ejemplo, postulación duplicada o convocatoria cerrada).
    \item \textbf{Reglas de permiso:} un caso por rol no autorizado, con respuesta 403.
\end{itemize}

Sobre tres casos de uso se aplicaron técnicas formales (Capítulo~\hyperref[cap:pruebas]{9}): caja blanca por caminos básicos, con la complejidad ciclomática de los métodos del \textit{backend} (6 para el envío de la postulación, 6 para la precalificación y 9 para la selección de finalistas); caja negra por partición de equivalencia y valores límite; y listas de comprobación para la verificación del equipo. Los tiempos de respuesta se diseñaron con los límites de los requisitos no funcionales.

\subsubsection{3.6.3. Implementar las pruebas}

Los casos de prueba se implementaron como pruebas automatizadas y como colecciones ejecutables. La Tabla~\ref{tab:pruebas_proyecto} muestra cuántas existen por parte del sistema: 406 archivos de prueba con 3\,548 pruebas escritas (\textit{backend}, panel web, aplicación móvil, módulo de IA y sitio de NEXWARE). Además:

\begin{itemize}
    \item Una colección de Bruno con 191 peticiones, una carpeta por caso de uso, para probar la API.
    \item Cinco pruebas de extremo a extremo con Playwright sobre el panel web.
    \item Un \textit{script} de k6 para la prueba de carga.
\end{itemize}

Las pruebas se ejecutan automáticamente en cada \textit{pull request} mediante GitHub Actions (sección 9.5.3).

\subsubsection{3.6.4. Realizar las pruebas de integración}

La integración se prueba con 22 clases de integración del \textit{backend}, que recorren el caso de uso completo (controlador, servicio y repositorio) contra una base PostgreSQL real levantada con Testcontainers. Esto comprueba también las migraciones de la base de datos y las restricciones reales del motor, que una base simulada no reproduce. La integración entre el \textit{backend} y el módulo de IA se cubre con pruebas que verifican que el cálculo del puntaje del módulo de Python coincide con el del \textit{backend}; en las pruebas del \textit{backend} el módulo de IA se reemplaza por simulaciones para controlar cada situación.

\subsubsection{3.6.5. Realizar las pruebas de sistema}

Las pruebas de sistema se hicieron con el \textit{backend} en perfil de desarrollo sobre una base recién creada:

\begin{itemize}
    \item \textbf{API (Bruno):} las 191 peticiones cubren los 26 casos de uso y las notificaciones, con su camino correcto, sus errores de negocio y los permisos por rol.
    \item \textbf{Panel web (Playwright):} cuatro flujos en un navegador real: inicio de sesión, registro y postulación hasta guardar el borrador, auditoría y adjudicación, y consulta de la bitácora.
    \item \textbf{Rendimiento:} la prueba de carga de k6 y la medición de los tiempos de respuesta del panel desplegado (sección 9.4).
    \item \textbf{Verificación manual:} la revisión de los tres casos de uso con listas de comprobación (sección 9.3.4).
\end{itemize}

\subsubsection{3.6.6. Evaluar las pruebas}

El resultado de la ejecución del 18 de noviembre de 2026 fue el siguiente:

\begin{itemize}
    \item \textbf{\textit{Backend}:} 2\,032 pruebas ejecutadas (una omitida), todas aprobadas, con 97,0\,\% de cobertura de líneas.
    \item \textbf{Panel web:} 1\,296 pruebas, todas aprobadas, con 83,4\,\% de cobertura de líneas.
    \item \textbf{Aplicación móvil:} 310 pruebas, todas aprobadas, con 78,7\,\% de cobertura de líneas.
    \item \textbf{API (Bruno):} 191 peticiones, 444 aserciones y 21 pruebas, sin ningún fallo.
    \item \textbf{Extremo a extremo (Playwright):} 5 de 5 pruebas aprobadas.
    \item \textbf{Carga (k6):} 0\,\% de peticiones fallidas y percentil 95 de 15,33 ms, por debajo de los umbrales de 1\,\% y 800 ms.
\end{itemize}

Los defectos que aparecieron se resolvieron durante las pruebas: una prueba de planillas esperaba un texto distinto al del PDF generado, y el \textit{script} de k6 usaba el rol equivocado para pedir los indicadores. Ambos fueron errores de las pruebas y no del sistema. El análisis estático con SonarQube Cloud calificó la seguridad del proyecto con C y registró 19 incidencias abiertas, que se revisan por severidad (sección 9.5.5).

Las limitaciones son que la cobertura no es del 100\,\%, que la colección de Bruno no cubre la subida real de archivos (CU-16) ni el enlace de recuperación por correo (CU-15), que no hay pruebas de extremo a extremo de la aplicación móvil y que la prueba de carga es ligera (20 usuarios durante un minuto). El detalle de los resultados está en el Capítulo~\hyperref[cap:pruebas]{9}.

\subsubsection{3.6.7. Pruebas de caja negra de los 26 casos de uso}

Las pruebas de caja negra verifican el comportamiento visible del sistema sin mirar el código: se parte de los datos de entrada y de los pasos que sigue el usuario, y se compara la respuesta con la esperada. La Tabla~\ref{tab:reporte_cu} presenta, para cada uno de los 26 casos de uso, los criterios de aceptación que se verifican, los datos de prueba, el procedimiento paso a paso y el resultado esperado. Cada fila se deriva de los criterios de aceptación de la especificación del caso de uso en el repositorio del sistema (carpeta de casos de uso), con un caso para el camino correcto, uno para cada error de negocio y uno para los permisos por rol; los códigos entre paréntesis (por ejemplo, CA03) remiten al criterio de origen. Los procedimientos se ejecutan con el sistema en perfil de desarrollo, que carga datos de prueba, y los cubre la colección de Bruno de la sección 9.5.5, que tiene una carpeta por caso de uso. Para tres de estos casos de uso (CU-02, CU-03 y CU-05), la caja negra se desarrolla con las técnicas de partición de equivalencia y de valores límite en la sección 9.3.2, y se complementa con la caja blanca de la sección 9.3.1.

{\footnotesize\linespread{1}\selectfont
\begin{longtable}{@{}>{\raggedright\arraybackslash}p{1.2cm}>{\raggedright\arraybackslash}p{3.7cm}>{\raggedright\arraybackslash}p{3.0cm}>{\raggedright\arraybackslash}p{4.4cm}>{\raggedright\arraybackslash}p{3.4cm}@{}}
\caption{Pruebas de caja negra de los 26 casos de uso}\label{tab:reporte_cu}\\
\toprule
\textbf{CU} & \textbf{Criterios de aceptación} & \textbf{Datos de prueba} & \textbf{Procedimiento de prueba} & \textbf{Resultado esperado} \\
\midrule
\endfirsthead
\toprule
\textbf{CU} & \textbf{Criterios de aceptación} & \textbf{Datos de prueba} & \textbf{Procedimiento de prueba} & \textbf{Resultado esperado} \\
\midrule
\endhead
\bottomrule
\endfoot
\textbf{CU-01} & - Registro con CI y correo nuevos crea una cuenta de Postulante y abre la sesión (CA01).\newline - CI o correo ya registrados responden 409 sin crear la cuenta (CA02, CA03).\newline - Cinco intentos fallidos bloquean la cuenta por 5 minutos (CA10). & Nombre, CI, registro universitario y correo no registrados; contraseña de 6 o más caracteres con su confirmación; cuenta existente del Administrador DUBSS. & 1. Abrir el registro, llenar el formulario y enviarlo.\newline 2. Cerrar sesión.\newline 3. Registrar otra cuenta con el mismo CI y luego con el mismo correo.\newline 4. Iniciar sesión con contraseña incorrecta cinco veces seguidas y una sexta con la correcta. & Se crea una sola cuenta con rol Postulante y se abre la sesión; el CI o el correo repetido responden 409; tras los cinco intentos la cuenta queda bloqueada aunque la contraseña sea correcta. \\
\midrule
\textbf{CU-02} & - Con formulario y documentos completos la postulación pasa a «pendiente» y devuelve el comprobante (CA03).\newline - Si falta un campo o documento obligatorio responde 400 ERR\_ENV\_04 (CA04).\newline - Un envío repetido devuelve el mismo comprobante sin duplicar (CA08). & Convocatoria abierta; ingresos familiares 2~500; 3 dependientes; procedencia Provincia; sin vivienda propia; hogar monoparental; promedio 72,5; documentos obligatorios adjuntos. & 1. Iniciar sesión como Postulante.\newline 2. Abrir la convocatoria y completar el formulario.\newline 3. Adjuntar los documentos y enviar.\newline 4. Repetir el envío.\newline 5. Enviar otra postulación con un campo obligatorio vacío. & La postulación queda «pendiente» con fecha de envío y código; el reintento devuelve el mismo comprobante; con el campo vacío responde 400 y el estado no cambia. \\
\midrule
\textbf{CU-03} & - Una postulación «pendiente» recibe un puntaje vigente de 0 a 70 con fecha y versión del modelo (CA01).\newline - Si la IA no responde en 5 s, la postulación sigue «pendiente» y se reintenta cada 5 minutos (CA02).\newline - Una solicitud repetida no genera un segundo puntaje (CA04). & Postulación «pendiente» con respuestas completas; módulo de IA disponible y luego detenido. & 1. Enviar la postulación (CU-02).\newline 2. Consultar su estado hasta que cambie.\newline 3. Repetir la solicitud de cálculo.\newline 4. Detener el módulo de IA y enviar otra postulación; volver a encenderlo. & La postulación pasa a «calificada» con un único puntaje entre 0 y 70; la solicitud repetida no crea otro; con la IA detenida queda «pendiente» y se califica sola al volver el servicio. \\
\midrule
\textbf{CU-04} & - Lista las postulaciones enviadas, de la más reciente a la más antigua, con gestión, fecha y estado (CA01).\newline - El detalle muestra el historial con fecha y hora, sin mostrar quién hizo el cambio (CA04).\newline - Una postulación «devuelta» muestra las observaciones y «Corregir y reenviar» (CA05). & Postulante con postulaciones en varios estados (calificada, devuelta); postulante sin postulaciones. & 1. Iniciar sesión como Postulante.\newline 2. Abrir «Mis postulaciones» y filtrar por gestión.\newline 3. Abrir el detalle de una postulación devuelta.\newline 4. Entrar con el postulante sin postulaciones. & La lista aparece ordenada con su estado; el detalle muestra el historial y las observaciones de la devolución; sin postulaciones se muestra un mensaje informativo. \\
\midrule
\textbf{CU-05} & - El panel muestra los cupos, los adjudicados y el puntaje mínimo (CA01).\newline - Se adjudica solo con auditoría «Validado», puntaje total igual o mayor al mínimo y cupos disponibles (CA04).\newline - No adjudicar sin motivo responde 400; sin auditoría, 409 ERR\_ADJ\_02 (CA06, CA07). & Convocatoria cerrada «[PRUEBA] E04» con 3 cupos; finalistas con y sin auditoría; motivo de hasta 500 caracteres. & 1. Iniciar sesión como Administrador DUBSS y abrir «Auditoría y adjudicación».\newline 2. Seleccionar los finalistas.\newline 3. Abrir un finalista auditado y adjudicar tras confirmar.\newline 4. Abrir otro y «No adjudicar» sin motivo; repetir con motivo.\newline 5. Intentar decidir sobre un finalista sin auditoría. & La postulación pasa a «adjudicada» y el resumen muestra un adjudicado y dos cupos disponibles; sin motivo responde 400; sin auditoría responde 409. \\
\midrule
\textbf{CU-06} & - Ordena las postulaciones calificadas de mayor a menor puntaje vigente (CA01).\newline - No incluye postulaciones en otros estados (CA03).\newline - Paginada: en la página 2 con 20 filas, la primera tiene la posición 21 (CA06). & Convocatoria «[PRUEBA] Volumen» con 42 postulaciones calificadas; Administrador DUBSS. & 1. Iniciar sesión como Administrador DUBSS.\newline 2. Abrir «Lista priorizada» y elegir la convocatoria.\newline 3. Pasar a la página 2 con 20 filas.\newline 4. Buscar por nombre o CI y filtrar por carrera.\newline 5. Exportar a PDF. & La lista aparece ordenada por puntaje; la página 2 empieza en la posición 21; los filtros conservan la posición de la lista completa; la exportación queda en la bitácora. \\
\midrule
\textbf{CU-07} & - Una auditoría «Validado» con entrega física recibida se guarda con resultado, auditor y fecha (CA03).\newline - «Observado» o «No válido» sin observaciones responde 400 (CA04).\newline - Sin la documentación física registrada responde 409 (CA07). & Finalistas de la convocatoria «[PRUEBA] E04», con y sin entrega física; puntaje académico a registrar. & 1. Abrir «Finalistas» y el expediente de un finalista.\newline 2. Intentar auditarlo sin la entrega registrada.\newline 3. Registrar la recepción de la documentación.\newline 4. Registrar la auditoría «Validado».\newline 5. Registrar otra «Observado» sin observaciones. & Sin entrega física responde 409; con ella la auditoría se guarda y se toma como vigente (las anteriores se conservan); «Observado» sin observaciones responde 400. \\
\midrule
\textbf{CU-08} & - Con resultados publicados genera el PDF y registra el acta con número (CA01).\newline - Sin resultados publicados responde 409 ERR\_ACT\_01; si ya existe acta, 409 ERR\_ACT\_02 (CA02, CA03).\newline - El número tiene el formato NNN/AAAA, correlativo y único (CA06). & Convocatoria «[PRUEBA]» con resultados publicados; otra sin publicar. & 1. Iniciar sesión como Administrador DUBSS y abrir «Actas».\newline 2. Elegir la convocatoria publicada y ver la vista previa.\newline 3. Generar el PDF y confirmar que es definitivo.\newline 4. Intentar generar otra acta de la misma convocatoria.\newline 5. Elegir la convocatoria sin publicar. & Se guarda un acta con su número y su PDF; la segunda y la de la convocatoria sin publicar responden 409. \\
\midrule
\textbf{CU-09} & - Con becas vigentes y un periodo válido genera la planilla apaisada y responde 201 (CA01, CA02).\newline - Fecha final anterior a la inicial responde 400 ERR\_PLA\_01; más de 31 días, 400 ERR\_PLA\_03 (CA03, CA05).\newline - Sin becas vigentes responde 409 ERR\_PLA\_04 (CA06). & Convocatoria con padrón de becados vigentes; periodos de 31 y de 32 días dentro de la gestión académica; convocatoria sin becas. & 1. Iniciar sesión como Administrador DUBSS y abrir «Planillas».\newline 2. Elegir la convocatoria e ingresar el periodo.\newline 3. Ver la vista previa y generar la planilla.\newline 4. Repetir con fecha final anterior y con 32 días.\newline 5. Probar con la convocatoria sin becas vigentes. & Se genera el PDF solo con los becados vigentes, ordenados por apellido y nombre; los periodos inválidos responden 400 y la convocatoria sin becas, 409. \\
\midrule
\textbf{CU-10} & - Una beca vigente sin asistencia hoy registra la asistencia y responde 201 (CA03).\newline - Si ya hay asistencia hoy responde 200 con yaRegistrada y no duplica (CA04).\newline - Beca suspendida o dada de baja responde 409 ERR\_ASI\_02; CI sin beca, 404 (CA06, CA02). & CI de un becado vigente, de uno suspendido y de un CI inexistente; cuenta del Personal del Comedor. & 1. Iniciar sesión como Personal del Comedor.\newline 2. Ingresar el CI del becado vigente y buscar.\newline 3. Registrar la asistencia.\newline 4. Repetir el registro.\newline 5. Buscar el CI del becado suspendido y el inexistente. & La primera vez se crea una asistencia con la hora; la repetición no crea otra; el becado suspendido responde 409 y el CI inexistente, 404. \\
\midrule
\textbf{CU-11} & - Muestra la planilla y los registros digitales de su periodo (CA01).\newline - Un cotejo consistente se guarda con usuario, fecha y resultado (CA04).\newline - Coincidencias y discrepancias que no suman el total responden 400 ERR\_COT\_02 (CA05). & Planilla generada (CU-09) con registros digitales de asistencia; cantidades de la planilla física firmada. & 1. Iniciar sesión como Administrador DUBSS y abrir «Cotejo de asistencia».\newline 2. Elegir la convocatoria y la planilla.\newline 3. Ingresar las coincidencias y discrepancias de la planilla física.\newline 4. Guardar; repetir con cantidades que no sumen el total. & El cotejo consistente se guarda con su resultado; el inconsistente responde 400 y no guarda nada. \\
\midrule
\textbf{CU-12} & - Muestra los indicadores de la gestión vigente con su fecha de actualización (CA01).\newline - Filtra por gestión o por convocatoria y actualiza los gráficos (CA02, CA03).\newline - Sin datos muestra un mensaje y ceros; los borradores no se cuentan (CA04, CA06). & Cuenta de Directivo; una gestión con datos y otra sin datos. & 1. Iniciar sesión como Directivo.\newline 2. Abrir «Dashboard de indicadores».\newline 3. Elegir otra gestión y luego una convocatoria.\newline 4. Elegir la gestión sin datos. & Se actualizan los cuatro indicadores y los gráficos con cada filtro; la gestión sin datos muestra los indicadores en cero y un mensaje, sin error. \\
\midrule
\textbf{CU-13} & - Crea la convocatoria en «BORRADOR» con sus documentos requeridos (CA01).\newline - Cierre igual o anterior a la apertura, o sin documento obligatorio, responde 400 (CA02, CA03).\newline - Abrir una convocatoria con otra «ABIERTA» en la misma gestión responde 409 (CA05). & Gestión académica vigente; fechas de apertura y cierre válidas; cupos, puntaje mínimo y finalistas en rango (valores de ejemplo); documentos requeridos. & 1. Iniciar sesión como Administrador DUBSS y abrir «Convocatorias».\newline 2. Crear una convocatoria con datos válidos.\newline 3. Intentar crear otra con el cierre anterior a la apertura.\newline 4. Abrirla; intentar abrir una segunda en la misma gestión. & La convocatoria queda en «BORRADOR» y pasa a «ABIERTA»; las fechas inválidas responden 400 y la segunda apertura, 409. \\
\midrule
\textbf{CU-14} & - Crea una cuenta activa del personal con el rol elegido (CA01).\newline - CI o correo repetidos responden 409; el rol Postulante responde 400 (CA02, CA03, CA05).\newline - La lista muestra solo cuentas del personal, paginadas y filtrables (CA06). & Nombre, CI y correo nuevos; roles Administrador DUBSS, Directivo y Personal del Comedor. & 1. Iniciar sesión como Administrador DUBSS y abrir la gestión de usuarios.\newline 2. Crear una cuenta con rol de personal.\newline 3. Intentar crearla con el rol Postulante y con un CI repetido.\newline 4. Cambiar el rol y desactivar y reactivar la cuenta. & La cuenta se crea activa y los cambios se reflejan en la lista; el rol Postulante responde 400 y el CI repetido, 409. \\
\midrule
\textbf{CU-15} & - Un correo de cuenta activa recibe un enlace de un solo uso (CA01).\newline - Un correo no registrado recibe el mismo mensaje neutro y no se envía correo (CA02).\newline - Un enlace vencido (más de 60 minutos) o ya usado se rechaza (CA04, CA05). & Correo de una cuenta activa y uno no registrado; contraseña nueva de 6 o más caracteres con su confirmación. & 1. Desde el inicio de sesión elegir «¿Olvidó su contraseña?».\newline 2. Solicitar el enlace con el correo registrado y con el no registrado.\newline 3. Abrir el enlace y definir la contraseña nueva.\newline 4. Abrir el mismo enlace por segunda vez. & Se envía un solo enlace y la contraseña cambia; el segundo uso se rechaza y no cambia nada. \\
\midrule
\textbf{CU-16} & - Un archivo válido de un tipo de la lista se guarda y se muestra adjuntado (CA02).\newline - Un formato no permitido o un archivo que excede el tamaño máximo responde 400 (CA03, CA04).\newline - Con la postulación ya enviada responde 409 (CA06). & Postulación en borrador de una convocatoria abierta; un archivo válido, uno de formato no permitido y uno muy grande. & 1. Iniciar sesión como Postulante y abrir la lista de documentos requeridos.\newline 2. Adjuntar el archivo válido.\newline 3. Intentar adjuntar el de formato no permitido y el muy grande.\newline 4. Enviar la postulación e intentar reemplazar un documento. & El archivo válido queda adjuntado; los otros responden 400 y no se guardan; tras el envío se responde 409. \\
\midrule
\textbf{CU-17} & - Guarda las respuestas escritas aunque falten campos y documentos (CA01).\newline - Un valor inválido responde 400 ERR\_ENV\_05 sin guardar ningún campo (CA03).\newline - Convocatoria cerrada responde 409 ERR\_VAL\_03; postulación ya enviada, 409 ERR\_ENV\_03 (CA04, CA05). & Postulación en borrador con algunos campos llenos; un valor fuera de rango; una postulación ya enviada. & 1. Iniciar sesión como Postulante y llenar parte del formulario.\newline 2. Elegir «Guardar borrador».\newline 3. Guardar con un valor fuera de rango.\newline 4. Intentar guardar sobre una postulación ya enviada. & El borrador se guarda con los campos vacíos como están y se recupera al volver; el valor inválido responde 400; la postulación enviada, 409. \\
\midrule
\textbf{CU-18} & - Un postulante elegible sin otra postulación vigente en la gestión es «válida» (CA01).\newline - Una postulación vigente en la misma gestión lo vuelve «no válida»; un borrador o una de otra gestión no cuentan (CA02, CA03, CA05).\newline - Un postulante no elegible recibe el motivo específico (CA06). & Postulante sin postulaciones; postulante con una postulación calificada en la misma gestión; postulante con un borrador. & 1. Enviar una postulación con el postulante sin postulaciones.\newline 2. Intentar enviar otra con el que ya tiene una vigente.\newline 3. Enviar con el que solo tiene un borrador. & La primera es válida; la segunda responde 409 ERR\_VAL\_01 (o 403 ERR\_VAL\_02 si no es elegible); el borrador no cuenta como duplicado. \\
\midrule
\textbf{CU-19} & - Devolver con motivo válido cambia la postulación a «devuelta» y lo anota en el historial (CA01).\newline - Motivo vacío o de más de 500 caracteres responde 400; si ya no está «calificada», 409 (CA02, CA03, CA04).\newline - Avisa al postulante por bandeja, push y correo (CA06). & Finalista con postulación «calificada»; motivo válido; motivo en blanco. & 1. Abrir el expediente de un finalista.\newline 2. Elegir «Devolver para corrección» e intentar confirmar con el motivo vacío.\newline 3. Confirmar con un motivo válido.\newline 4. Repetir la devolución sobre la misma postulación. & La postulación pasa a «devuelta», con historial y un solo aviso al postulante; el motivo vacío responde 400 y la repetición, 409. \\
\midrule
\textbf{CU-20} & - Archivar con motivo deja la postulación «archivada»; rechazar la deja «no adjudicada» (CA01, CA02).\newline - Motivo vacío o de más de 500 caracteres responde 400 (CA03, CA04).\newline - Resultados ya publicados o postulación que ya no está «calificada» responden 409 (CA05, CA06). & Finalistas con postulación «calificada»; motivo válido y motivo en blanco; convocatoria con resultados publicados. & 1. Abrir el expediente de un finalista.\newline 2. Archivar con motivo.\newline 3. Abrir otro y rechazar con motivo.\newline 4. Intentar archivar con el motivo vacío.\newline 5. Intentar en la convocatoria con resultados publicados. & Los estados cambian a «archivada» y «no adjudicada» con su historial; el motivo vacío responde 400 y la convocatoria publicada, 409. \\
\midrule
\textbf{CU-21} & - Lista los finalistas con el estado de su notificación y de su entrega física (CA01).\newline - Con plazo, lugar, horario y documentos guarda un aviso en la bandeja de cada uno y envía push y correo (CA02).\newline - Faltan datos o el plazo no es posterior a hoy: responde 400 (CA05, CA06). & Convocatoria con finalistas, uno de ellos sin teléfono registrado; plazo posterior a hoy; lugar, horario y documentos a presentar. & 1. Abrir «Finalistas» y elegir la convocatoria.\newline 2. Seleccionar los finalistas.\newline 3. Completar plazo, lugar, horario y documentos y enviar.\newline 4. Intentar enviar con el plazo de hoy y sin documentos. & Cada finalista recibe el aviso (el que no tiene teléfono, en la bandeja y por correo); los datos faltantes o el plazo inválido responden 400. \\
\midrule
\textbf{CU-22} & - Sin finalistas en «calificada» registra la publicación con usuario y fecha (CA01).\newline - Si quedan finalistas sin decidir responde 409 e indica cuántos faltan (CA02).\newline - Una publicación repetida responde 409 y avisa una sola vez a cada finalista (CA03, CA04, CA05). & Convocatoria con todas las decisiones tomadas y otra con finalistas sin decidir. & 1. Abrir «Publicar resultados» y elegir la convocatoria con decisiones pendientes.\newline 2. Intentar publicar.\newline 3. Elegir la convocatoria con todas las decisiones y publicar.\newline 4. Publicar otra vez. & La primera responde 409 con la cantidad pendiente; la segunda se publica y avisa a cada finalista; la repetición responde 409. \\
\midrule
\textbf{CU-23} & - Tras la publicación crea una beca vigente por cada adjudicado, sin duplicar (CA00).\newline - La baja o suspensión con motivo cambia el estado y se registra; sin motivo responde 400 (CA02, CA03).\newline - Una beca dada de baja se reemplaza con el siguiente postulante priorizado (CA05). & Convocatoria con resultados publicados y adjudicados sin alta; motivo de baja. & 1. Abrir «Padrón de becados» y elegir la convocatoria.\newline 2. Dar de alta a los adjudicados y repetir la operación.\newline 3. Suspender una beca sin motivo y luego con motivo.\newline 4. Dar de baja otra y reemplazarla. & Se crea una beca vigente por adjudicado, sin duplicados; la suspensión sin motivo responde 400; el reemplazo crea la beca del siguiente postulante. \\
\midrule
\textbf{CU-24} & - Muestra el código, el estado, la convocatoria y las fechas de la beca (CA01).\newline - Un usuario sin beca recibe 404 ERR\_MBE\_01 y se lo lleva a «Mis postulaciones» (CA03).\newline - Lista las asistencias paginadas y las filtra por rango de fechas (CA04, CA05). & Becado con asistencias registradas; becado con beca suspendida; postulante sin beca. & 1. Iniciar sesión como becado y abrir «Mi beca».\newline 2. Revisar el historial de asistencias y filtrar por un rango de fechas.\newline 3. Entrar con el becado suspendido.\newline 4. Entrar con un postulante sin beca. & Se ve el estado y el historial; la beca suspendida muestra el motivo; sin beca responde 404 y la pantalla lleva a «Mis postulaciones». \\
\midrule
\textbf{CU-25} & - Lista las acciones de la más reciente a la más antigua, paginadas (CA01).\newline - Filtra por usuario, acción o rango de fechas; «desde» posterior a «hasta» responde 400 ERR\_BIT\_02 (CA02, CA03).\newline - Un usuario que no es Administrador recibe 403; no hay operaciones para modificar o eliminar (CA04, CA05). & Acciones previas en el sistema (inicios de sesión, adjudicaciones); cuenta de Administrador y de otro rol. & 1. Iniciar sesión como Administrador DUBSS y abrir «Bitácora».\newline 2. Filtrar por usuario y por rango de fechas.\newline 3. Poner «desde» posterior a «hasta».\newline 4. Entrar a la ruta con una cuenta de otro rol. & La bitácora muestra solo las acciones que cumplen los filtros; las fechas invertidas responden 400 y el otro rol, 403. \\
\midrule
\textbf{CU-26} & - Muestra la versión vigente, sus ponderaciones y el historial de versiones (CA01).\newline - Ponderaciones que no suman 100 responden 400 ERR\_MOD\_01; las válidas crean la versión siguiente y archivan la anterior (CA02, CA03).\newline - Reprocesar recalcula las postulaciones calificadas; con resultados publicados responde 409 ERR\_MOD\_04 (CA05, CA06). & Ponderaciones que suman 100 y otras que no; convocatoria con postulaciones calificadas y otra con resultados publicados. & 1. Abrir «Parámetros del modelo».\newline 2. Guardar ponderaciones que no suman 100.\newline 3. Guardar ponderaciones válidas.\newline 4. Confirmar el reprocesamiento de una convocatoria.\newline 5. Intentar reprocesar la convocatoria publicada. & Las ponderaciones inválidas responden 400; las válidas crean la versión siguiente; el reprocesamiento recalcula los puntajes y la convocatoria publicada responde 409. \\
\bottomrule
\end{longtable}
}

\noindent\footnotesize\textit{Nota.} Elaboración propia a partir de los criterios de aceptación de los casos de uso. Los resultados esperados corresponden a la corrida de Bruno del 18 de noviembre de 2026 (191 peticiones sin fallos). La subida real de archivos del CU-16 y el enlace de recuperación por correo del CU-15 no se cubren en esa corrida (sección 9.5.5.1). Los valores de convocatoria, cupos y puntaje mínimo de los datos de prueba son de ejemplo; los definitivos están por definir con la DUBSS.\normalsize
